\chapter{Diskussion}\label{ch:diskussion}

Die Diskussion interpretiert die Ergebnisse inhaltlich und ordnet sie in den wissenschaftlichen Kontext ein.
Ziel ist es, die Bedeutung und Grenzen der Arbeit herauszuarbeiten.

Typische Punkte:
\begin{itemize}
\item Interpretation der Ergebnisse – warum treten beobachtete Effekte auf?
\item Bewertung im Vergleich zu anderen Studien oder Verfahren.
\item Grenzen der Methode: Datengröße, Annahmen, Messfehler, Generalisierbarkeit.
\item Praktische Relevanz der Resultate – was bedeuten sie für Anwender oder Forschung?
\end{itemize}

Ein kritischer, aber sachlicher Ton ist hier angebracht. Zeigen Sie sowohl Stärken als auch Schwächen Ihres Ansatzes auf.


